\begin{table}[H]
  \centering
  \caption{Perfil del hogar vulnerable (``entra en pobreza'', 2010$\to$2013) desagregado por trayectoria a 2016: transitorio ($n=355$, vuelve a salir) vs.\ persistente ($n=255$, se queda pobre). Valores ponderados por grupo; filas superiores, resumen agregado sobre las 53 variables; filas inferiores, ejemplos puntuales donde el promedio agregado (Sección~\ref{subsec:caracterizacion_grupos}) ocultaba una divergencia entre subgrupos.}
  \label{tab:perfil_split_trayectoria}
  \footnotesize
  \setlength{\tabcolsep}{4pt}
  \begin{tabular}{lrrrrr}
    \toprule
    \textbf{Variable} & \textbf{Siempre} & \textbf{Sale} & \textbf{Transitorio} & \textbf{Persistente} & \textbf{Nunca} \\
    \midrule
    Más cerca de ``sale'' que de ``nunca'' (de 53) & \multicolumn{2}{c}{--} & 39/53 & 44/53 & -- \\
    Más cerca de ``siempre'' que de ``sale'' (de 53) & \multicolumn{2}{c}{--} & 2/53 & 11/53 & -- \\
    \addlinespace
    Jefe asalariado & 34.4\% & 52.2\% & 64.9\% & 61.7\% & 57.9\% \\
    Jefe cotiza a pensión & 6.4\% & 21.3\% & 33.0\% & 12.0\% & 42.8\% \\
    Tiene deuda informal & 31.6\% & 22.7\% & 17.3\% & 41.9\% & 9.4\% \\
    Tiene deuda formal & 65.2\% & 74.6\% & 81.4\% & 57.5\% & 88.3\% \\
    Comunidad con transporte público & 54.6\% & 75.4\% & 84.9\% & 78.2\% & 78.2\% \\
    \bottomrule
  \end{tabular}
  \begin{minipage}{0.95\textwidth}
    \vspace{4pt}
    \footnotesize \textit{Nota:} de los 723 hogares que entran en pobreza entre 2010 y 2013, 610 tienen dato de la ola 2016 (113 se pierden por atrición, misma tasa que el resto del panel); de esos 610, 355 (58\%) vuelven a salir de la pobreza para 2016 y 255 (42\%) se quedan pobres. Fuente: cálculos propios, \texttt{perfil\_split\_transitorio\_persistente\_monetaria.csv} (\texttt{src/02\_build/eda\_perfil\_split\_trayectoria.py}).
  \end{minipage}
\end{table}